\section{Seguridad aplicativa}

\subsection{Descripcion}

Este documento define los lineamientos de seguridad aplicativa del proyecto alineados con \textit{Secure by Design}, buenas practicas OWASP y principios de produccion ya aprobados para Quiniela.

\subsection{Alcance}

Aplica a autenticacion, autorizacion, validacion de entradas, proteccion de configuracion, hardening base y manejo seguro de errores. No sustituye una auditoria de seguridad ni una politica de infraestructura detallada.

\subsection{Definiciones}

\begin{itemize}
  \item \textbf{Secure by Design}: seguridad incorporada desde el diseno y no agregada al final.
  \item \textbf{Superficie de ataque}: puntos donde el sistema recibe datos, ejecuta acciones o expone informacion.
  \item \textbf{Hardening}: medidas de reduccion de riesgo en configuracion y comportamiento del sistema.
\end{itemize}

\subsection{Reglas}

\begin{itemize}
  \item Todo input externo debe validarse de forma estricta en la frontera de entrada.
  \item Toda accion critica debe exigir autenticacion valida, autorizacion acorde al rol y trazabilidad suficiente.
  \item La configuracion sensible debe vivir fuera del codigo y gestionarse por entorno.
  \item Los errores expuestos al cliente no deben filtrar detalles internos innecesarios.
  \item La seguridad no puede debilitar reglas del dominio; debe reforzarlas.
\end{itemize}

\subsubsection{Controles minimos esperados}

\begin{itemize}
  \item validacion estricta de payloads y parametros;
  \item manejo seguro de sesion o tokens;
  \item proteccion de rutas administrativas;
  \item headers y hardening HTTP apropiados para produccion;
  \item politicas de contrasena y recuperacion de acceso controladas;
  \item proteccion contra abuso basico en autenticacion y acciones sensibles;
  \item sanitizacion y manejo seguro de archivos o evidencias cargadas si aplica.
\end{itemize}

\subsubsection{Autenticacion y autorizacion}

\begin{itemize}
  \item La autenticacion debe identificar de forma fiable al actor antes de permitir operaciones protegidas.
  \item La autorizacion debe evaluarse por rol y por estado funcional cuando la accion lo requiera.
  \item Una cuenta autenticada pero no \texttt{ACTIVA} no puede ejecutar acciones competitivas reservadas a jugadores habilitados.
  \item Las operaciones administrativas no deben confiar en ocultamiento de UI; su proteccion principal vive en backend.
\end{itemize}

\subsubsection{Validacion y manejo de entrada}

\begin{itemize}
  \item Los contratos de entrada deben definirse de forma explicita y rechazarse si no cumplen formato, rango o contexto funcional.
  \item El rango de goles permitido para picks debe respetar la regla oficial del producto: enteros entre \texttt{0} y \texttt{9}.
  \item Una entrada formalmente valida puede seguir siendo rechazada por estado del dominio, por ejemplo lock vencido o cuenta no activa.
  \item Las validaciones del frontend son complementarias; el backend es la autoridad final.
\end{itemize}

\subsubsection{Configuracion y secretos}

\begin{itemize}
  \item Ningun secreto debe quedar hardcodeado en codigo fuente ni en documentos de ejemplo.
  \item La configuracion operativa debe separarse por entorno y mantenerse fuera del repositorio cuando sea sensible.
  \item Los valores por defecto del sistema no deben abrir superficies inseguras en produccion.
\end{itemize}

\subsubsection{Manejo de errores}

\begin{itemize}
  \item Los errores funcionales deben informar rechazo de forma comprensible sin exponer detalle tecnico sensible.
  \item Los errores tecnicos deben registrarse con suficiente contexto interno para diagnostico.
  \item Una respuesta de error no debe filtrar internals de base de datos, stack trace ni configuracion.
\end{itemize}

\subsection{Edge Cases}

\begin{itemize}
  \item Si el frontend intenta una accion permitida por interfaz pero rechazada por backend, el sistema debe preferir seguridad y reconciliar la vista.
  \item Si una carga de comprobante o archivo existe en v1, su validacion y almacenamiento deben considerarse superficie sensible.
  \item Si una accion administrativa modifica resultados o scoring, la proteccion debe cubrir tanto acceso como auditoria posterior.
\end{itemize}

\subsection{Invariantes}

\begin{itemize}
  \item Ninguna accion critica opera sin autenticacion, autorizacion y trazabilidad acordes.
  \item Ningun input externo cruza la frontera del sistema sin validacion.
  \item Ninguna respuesta publica expone internals innecesarios del sistema.
  \item La configuracion sensible permanece fuera del codigo fuente.
\end{itemize}

\subsection{Preguntas abiertas}

\begin{itemize}
  \item Confirmar el checklist final de controles obligatorios de v1, incluyendo politica de rate limiting, manejo de archivos y nivel exacto de protecciones HTTP.
\end{itemize}

\subsection{Decisiones pendientes}

\begin{itemize}
  \item \texttt{Decision pendiente} \texttt{Bloquea backend}: cerrar el set minimo obligatorio de controles de autenticacion, gestion de sesion y proteccion contra abuso para v1.
  \item \texttt{Decision pendiente} \texttt{No bloqueante}: definir si el proyecto documentara una checklist OWASP interna separada o si estos controles vivirán solo en este documento.
\end{itemize}
